% Part: first-order-logic
% Chapter: sequent-calculus
% Section: rules-and-proofs

\documentclass[../../../include/open-logic-section]{subfiles}

\begin{document}

\iftag{FOL}
      {\olfileid{fol}{seq}{rul}}
      {\olfileid{pl}{seq}{rul}}

\olsection{નિયમો અને \usetoken{P}{derivation}}

નીચેના માટે, $\Gamma, \Delta, \Pi, \Lambda$ એ !!{sentence}sની સાન્ત
શ્રેણીઓ દર્શાવે તેમ માનો.

\begin{defn}[સિક્વન્ટ]
\emph{સિક્વન્ટ} એ નીચેના સ્વરૂપની અભિવ્યક્તિ છે:
\[
\Gamma \Sequent \Delta
\]
જ્યાં $\Gamma$ અને $\Delta$ એ ભાષા~$\Lang L$નાં !!{sentence}sની સાન્ત
(સંભવતઃ ખાલી) શ્રેણીઓ છે. $\Gamma$ને \emph{પૂર્વાંગ} અને $\Delta$ને
\emph{ઉત્તરાંગ} કહે છે.
\end{defn}

\begin{explain}
સિક્વન્ટ પાછળનો સહજ વિચાર આ છે: પૂર્વાંગનાં બધાં !!{sentence}s સાચાં હોય,
તો ઉત્તરાંગનાં !!{sentence}sમાંથી ઓછામાં ઓછું એક સાચું હોય. એટલે કે, જો
$\Gamma = \tuple{!A_1, \dots, !A_m}$ અને
$\Delta = \tuple{!B_1, \dots, !B_n}$ હોય, તો
$\Gamma \Sequent \Delta$ ત્યારે અને માત્ર ત્યારે જ સધાય જ્યારે
\[
(!A_1 \land \cdots \land !A_m) \lif (!B_1 \lor \cdots \lor
!B_n)
\]
સધાય. બે ખાસ કિસ્સા છે: $\Gamma$ ખાલી હોય તે અને $\Delta$ ખાલી હોય તે.
$\Gamma$ ખાલી હોય, એટલે કે $m = 0$ હોય, ત્યારે $\quad \Sequent \Delta$
ત્યારે અને માત્ર ત્યારે જ સધાય જ્યારે $!B_1 \lor \dots \lor !B_n$
સધાય. $\Delta$ ખાલી હોય, એટલે કે $n = 0$ હોય, ત્યારે
$\Gamma \Sequent \quad$ ત્યારે અને માત્ર ત્યારે જ સધાય જ્યારે
$\lnot(!A_1 \land \dots \land !A_m)$ સધાય. અનુરૂપ !!{sentence}
પ્રામાણ્ય હોય ત્યારે અને માત્ર ત્યારે જ સિક્વન્ટને પ્રામાણ્ય કહીએ છીએ.
\end{explain}

જો $\Gamma$ એ !!{sentence}sની શ્રેણી હોય, તો $\Gamma$ના જમણા છેડે
$!A$ જોડવાથી મળતા પરિણામ માટે $\Gamma, !A$ લખીએ છીએ (અને $\Gamma$ના
ડાબા છેડે $!A$ જોડવાથી મળતા પરિણામ માટે $!A, \Gamma$ લખીએ છીએ). જો
$\Delta$ પણ !!{sentence}sની શ્રેણી હોય, તો $\Gamma, \Delta$ એ બંને
શ્રેણીઓનું સળંગ જોડાણ છે.

\begin{defn}[આરંભિક સિક્વન્ટ]
\emph{આરંભિક સિક્વન્ટ} એ સિક્વન્ટ
\iftag{prvFalse,prvTrue}{નીચે આપેલાં સ્વરૂપોમાંથી કોઈ એક સ્વરૂપનું હોય છે:
  \begin{enumerate}
    \item $!A \Sequent !A$
    \tagitem{prvTrue}{$\quad \Sequent \ltrue$}{}
    \tagitem{prvFalse}{$\lfalse \Sequent \quad$}{}
  \end{enumerate}}
{$!A \Sequent !A$ સ્વરૂપનું હોય છે}, જ્યાં $!A$ ભાષાનું કોઈપણ
!!{sentence} છે.
\end{defn}

સિક્વન્ટ કલનમાં !!^{derivation}s એ સિક્વન્ટોનાં અમુક વૃક્ષો છે. તેમાં
સૌથી ઉપરના સિક્વન્ટો આરંભિક સિક્વન્ટો હોય છે; અને કોઈ સિક્વન્ટ એક કે બે
બીજા સિક્વન્ટોની નીચે હોય, તો તે અનુમાનના કોઈ નિયમ વડે તેમાંથી યોગ્ય
રીતે મળવો જોઈએ. $\Log{LK}$ના નિયમો બે મુખ્ય પ્રકારમાં વહેંચાય છે:
\emph{તાર્કિક} નિયમો અને \emph{સંરચનાત્મક} નિયમો. તાર્કિક નિયમોનાં નામ
નીચલા સિક્વન્ટમાં $!A$ અને/અથવા $!B$ ધરાવતા !!{sentence}ના
!!{main
  operator} પરથી પાડવામાં આવે છે. દરેક તાર્કિક નિયમનાં બે રૂપ હોય
છે: એક એવું રૂપ જે ડાબી બાજુ !!{operator} ધરાવતું !!{sentence} હોય એવા
સિક્વન્ટનું અનુમાન કરે, અને બીજું એવું રૂપ જે જમણી બાજુ તે !!{sentence}
હોય એવા સિક્વન્ટનું અનુમાન કરે.

\end{document}
